% =====================================================
% ROTEIRO PRÁTICO 01
% =====================================================

\section{Roteiro Prático 01 — Diagnóstico e Retomada do Projeto}

% =====================================================
% OBJETIVOS
% =====================================================

\begin{boxObjetivos}
\textbf{Objetivos da Aula}

Ao final desta aula o estudante será capaz de:

\begin{itemize}
\item Recuperar e organizar materiais do projeto anterior;
\item Identificar o estágio atual do projeto;
\item Analisar criticamente resultados obtidos;
\item Reconhecer limitações técnicas e metodológicas;
\item Estruturar um diagnóstico inicial formal do projeto.
\end{itemize}
\end{boxObjetivos}

% =====================================================
% INTRODUÇÃO
% =====================================================

\subsection{Introdução}

O Projeto Integrador II não inicia um novo projeto do zero.

Ele parte da maturação do projeto anteriormente desenvolvido, promovendo:

\begin{itemize}
\item Revisão crítica dos resultados;
\item Organização técnica dos artefatos;
\item Replanejamento estruturado;
\item Evolução técnica com base em diagnóstico;
\item Consolidação científica da solução.
\end{itemize}

Antes de avançar, é fundamental compreender com clareza o estágio atual do projeto.

Esta aula é dedicada à organização, análise e reflexão técnica.

% =====================================================
% PARTE 1
% =====================================================

\subsection{Parte 1 — Organização e Recuperação do Projeto (25 minutos)}

\begin{boxAtividade}[Atividade 1 — Recuperação dos Materiais]

Cada grupo deverá:

\begin{itemize}
\item Localizar o repositório do projeto anterior;
\item Verificar se o código está acessível e organizado;
\item Identificar se há datasets utilizados;
\item Reunir relatórios, apresentações e documentos produzidos;
\item Listar funcionalidades implementadas;
\item Listar funcionalidades previstas que não foram concluídas.
\end{itemize}

Ao final, elaborar uma lista estruturada contendo:

\begin{itemize}
\item O que existe;
\item O que está incompleto;
\item O que não funciona;
\item O que precisa ser reorganizado.
\end{itemize}

\end{boxAtividade}

% =====================================================
% PARTE 2
% =====================================================

\subsection{Parte 2 — Apresentação Sintética do Projeto (5 minutos por grupo)}

\begin{boxAtividade}[Atividade 2 — Pitch Técnico]

Cada grupo deverá apresentar:

\begin{itemize}
\item Problema abordado no PI-I;
\item Objetivos definidos originalmente;
\item Cronograma previsto;
\item Resultados alcançados;
\item Principais dificuldades encontradas;
\item Limitações atuais do projeto.
\end{itemize}

A apresentação deve ser objetiva e técnica.

\end{boxAtividade}

% =====================================================
% PARTE 3
% =====================================================

\subsection{Parte 3 — Diagnóstico Estruturado}

\begin{boxAtividade}[Atividade 3 — Documento de Diagnóstico]

Responder às seguintes questões no caderno de notas do projeto:

\begin{enumerate}
\item O projeto está funcional? Em qual nível?
\item O problema original ainda é relevante?
\item O escopo estava adequado ou superdimensionado (tentamos abraçar o mundo)?
\item Quais foram os principais erros de planejamento?
\item O grupo permanece o mesmo?
\item Há dados suficientes para experimentação?
\item Quais melhorias são prioritárias?
\item O projeto pode gerar 2 ou 3 micro-experimentos mensuráveis?
\end{enumerate}

O diagnóstico deverá ser redigido de forma técnica e objetiva (1 a 2 páginas) e deverá compor uma seção específica no carderno de notas do projeto chamada Diagnóstico do PI-1.

\end{boxAtividade}


% =====================================================
% PARTE 4
% =====================================================

\subsection{Parte 4 — Criação do Relatório Técnico no Overleaf}

\begin{boxAtividade}[Atividade 4 — Iniciando o Relatório Oficial do Projeto]

Nesta etapa, cada grupo deverá criar o documento formal do projeto utilizando o template institucional IFSC disponibilizado pelo professor.

\textbf{Etapa 1 — Baixar o Template no GitHub}

\begin{enumerate}
    \item Acessar o repositório oficial do template:
    
    \texttt{https://github.com/chameoandre/ifsc-gpb-relatorio-template}
    
    \item Clicar no botão verde \textbf{Code};
    \item Selecionar a opção \textbf{Download ZIP};
    \item Salvar o arquivo no computador;
    \item Extrair (descompactar) o arquivo ZIP.
\end{enumerate}

\vspace{0.4cm}

\textbf{Etapa 2 — Criar Projeto no Overleaf}

\begin{enumerate}
    \item Acessar \texttt{https://www.overleaf.com};
    \item Criar conta (caso não possua);
    \item Clicar em \textbf{New Project};
    \item Selecionar \textbf{Upload Project};
    \item Enviar o conteúdo da pasta extraída do template;
    \item Confirmar criação do projeto.
\end{enumerate}
\end{boxAtividade}

\vspace{0.4cm}


\newpage
\begin{boxAtividade}[Atividade 4 — Iniciando o Relatório Oficial do Projeto Etapa 3]

\textbf{Etapa 3 — Primeira Edição Obrigatória}

Após abrir o projeto no Overleaf:

\begin{itemize}
    \item Atualizar o bloco de identificação do documento;
    \item Inserir o título provisório do projeto;
    \item Preencher integralmente a seção:
    \textbf{Diagnóstico Inicial do Projeto};
    \item Compilar o documento e verificar se não há erros.
\end{itemize}

\textbf{Observação:} Este documento será atualizado ao longo de todo o semestre e servirá como base para a entrega final.

\end{boxAtividade}



\newpage

% =====================================================
% CHECKLIST
% =====================================================

\subsection{Checklist Final da Aula}

\begin{boxAtividade}[Verificação Técnica]

Antes de finalizar a aula, verifique:

\begin{itemize}
\item O repositório está acessível e organizado;
\item O grupo identificou claramente o estágio atual do projeto;
\item O diagnóstico foi redigido de forma estruturada;
\item O caderno de notas foi criado;
\item Há uma visão preliminar de possíveis micro-experimentos.
\end{itemize}

\end{boxAtividade}

% =====================================================
% REFLEXÃO FINAL
% =====================================================

\subsection*{Reflexão Técnica}

Projetos não evoluem apenas por implementação de código.

Eles evoluem por análise crítica, organização, planejamento e validação estruturada.

O Projeto Integrador II exige maturidade técnica.

O diagnóstico realizado hoje é a base para todas as etapas seguintes.
